\section{Conclusion}
We compared eight document-level detectors for RAG poisoning under one protocol
and asked where they fail rather than how high one number can go. The best method,
a logistic-regression fusion of the three original signals and a sentence
classifier, flags $\resTPRFiveLRsen$ of held-out attacks at a 5\% false-alarm
budget, and $\resTPROneLRsen$ at 1\%.
That recall is uneven: it holds up on attacks that carry an instruction, but at the
1\% budget only about one in ten of the passages that carry none is caught. A control
that sees only character $n$-grams shows that part of the classifier's recall on
those passages reflects corpus style, and when a whole source is held out the pooled
threshold either flags nearly every clean document or almost no attacks. Because
attacks are rare, an alert is often a false alarm: at one poisoned document in a
thousand the fused model's precision is only $\resPrecFiveLRsenPrevTenth$.
On cost, perplexity accounts for most of the detector time, and a cascade that
stops early on clearly clean documents saves more than half of the expected time on
clean traffic, at a cost of 0.08 recall at the 5\% budget and 0.01 at 1\%
(Table~\ref{tab:cascade}).

These results come from benchmark attacks, about $\resNNegA$ clean documents per
test role, and three overlapping splits. We did not test obfuscation or adaptive
attackers, per-source or conformal thresholds, or the effect of filtering on
downstream answers (Section~\ref{sec:limitations}). The practical lesson is that a
cheap detector should be calibrated on the source it will protect, evaluated on
documents from that source it has not seen, and judged by the precision of its
alerts, not by recall alone.